% !TEX root = ../main.tex
\begin{table}[htbp]
  \centering
  \small
  \caption{Claves de deduplicación consideradas (ADR-0001).}
  \label{tab:claves_dedup}
  \begin{tabularx}{\textwidth}{@{}>{\raggedright\arraybackslash}p{3.1cm} c c Y@{}}
    \toprule
    \textbf{Clave}      & \textbf{Conjuntos} & \textbf{Copias} & \textbf{Comportamiento observado}                      \\
    \midrule
    Solo contenido      & 41                 & 81              & 
    Agrupa las tres parejas dudosas, incluida la que no debe agruparse. Descartada.                                     \\
    \addlinespace
    (nombre, contenido) & 38                 & 78              & 
    Evita la fusión errónea a costa de dejar sin agrupar dos parejas que sí eran la misma asignatura. \textbf{Elegida}. \\
    \bottomrule
  \end{tabularx}
  \begin{minipage}{0.95\textwidth}
    \vspace{4pt}
    \tiny
    \textcolor{borrador2}{\textbf{Notas:} \textbf{Conjuntos}: grupos de dos o más guías con la misma clave, que se indexan como una sola unidad; \textbf{Copias}: guías que sobran al quedarse una por grupo (con la clave elegida, 116 guías en 38 grupos). Fuente: Elaboración propia.}
  \end{minipage}
\end{table}
